\section{Full enclosing cuts and their minimizing frontiers}
\label{ct:section}
\label{sec:cuts}

The area in the theorem is an intrinsic exterior invariant. In particular,
a portion of a cut coinciding with the original boundary still contributes
area. We establish the compactness and variation properties of this
invariant without using the constraint equations.

Write $a_g(S)=A_{\min}(S)$, with the convention of
Definition~\ref{def:full-cut}. All surface orientations point toward the
end. This normal is the inward normal of the exterior region bounded
by the cut.

\begin{lemma}[Flux through the entire intrinsic boundary]
\label{ct:flux}
If $\beta$ is a smooth closed two-form on $\Omega$, smooth up to $S$, then
for every full enclosing cut $\Gamma$ and every sufficiently distant
coordinate sphere $S_r$,
\begin{equation}
\label{ct:flux-identity}
 \int_\Gamma\beta=\int_{S_r}\beta=\int_S\beta.
\end{equation}
No extension of $\beta$ across $S$ is required.
\end{lemma}

\begin{proof}
Let $D$ be the exterior region defining $\Gamma$. Choose $r$ so that
$\Gamma$ lies strictly inside $S_r$ and a neighborhood of $S_r$ lies in
the manifold interior of $D$. The truncated region $D_r$ is a compact
smooth manifold with intrinsic boundary $\Gamma\sqcup S_r$. Its boundary
orientation on $\Gamma$ is opposite to the chosen end orientation.
Stokes' Theorem gives
$0=\int_{D_r}d\beta=\int_{S_r}\beta-\int_\Gamma\beta$.
The same argument on the truncation of $\Omega$ proves the other
equality. Contact between $\Gamma$ and $S$ occurs on the smooth intrinsic
boundary of $D_r$; it neither deletes a boundary term nor creates another
face. Every boundary component is included.
\end{proof}

For the two closed flux forms $\alpha_1,\alpha_2$ in
Equation~\ref{eq:flux-forms}, Equation~\ref{ct:flux-identity} gives the two
conserved charges individually.

\begin{proposition}[Strict positivity of the full-cut area]
\label{ct:full-cut-positive}
Under the geometric hypotheses of Definition~\ref{def:data},
\begin{equation}
\label{ct:positive-area}
 0<a_g(S)\leq\operatorname{Area}_g(S)<\infty.
\end{equation}
There is, more precisely, a smooth closed two-form $\beta$ on $\Omega$
with finite positive comass bound $C_\beta$ and
$\int_\Gamma\beta=1$ for every full cut. Thus $a_g(S)\geq C_\beta^{-1}$.
\end{proposition}

\begin{proof}
Choose a point of $S$. Its boundary collar, connectedness of $\Omega$,
and the single coordinate end give an embedded proper ray starting
transversely at that point and eventually equal to a straight coordinate
ray. Indeed, join the collar to the end by a path, perturb its compact
part to an embedded arc in dimension three, shorten at its last crossing
of a large coordinate sphere, and attach the outward straight ray.
The two joins can be smoothed in disjoint coordinate balls.

The normal disk bundle of the ray is trivial because its base is an
interval. The tubular-neighborhood construction on its compact portion,
followed by the straight tail, gives an embedded tube
\[
 \Phi:[0,\infty)\times\mathbb D^2\longrightarrow\Omega
\]
whose initial disk lies in $S$, whose positive-parameter part lies in
$\operatorname{int}\Omega$, and whose tail has a fixed positive coordinate
radius. Let $\omega$ be a smooth two-form compactly supported in the
interior of the transverse disk, oriented so that
$\int_{\mathbb D^2}\omega=1$. Define $\beta$ on the tube by
$\Phi^*\beta=\operatorname{pr}_{\mathbb D^2}^*\omega$, and extend by zero.
The transverse compact support makes this extension smooth at the lateral
boundary; it is smooth up to $S$ and closed.

Its comass is bounded on the compact portion and on the straight tail,
the latter because the coordinate radius is fixed and
$g_{ij}\to\delta_{ij}$. Thus $C_\beta=\sup_\Omega\|\beta\|_{*,g}<\infty$.
A sufficiently distant sphere meets the supporting tail in a positively
oriented graph over the entire support of $\omega$ and has flux one.
Lemma~\ref{ct:flux} gives the same flux on every full cut. Hence
\[
 1=\left|\int_\Gamma\beta\right|
 \leq C_\beta\operatorname{Area}_g(\Gamma).
\]
Finally $S$ itself is an admissible cut.
\end{proof}

\subsection{An auxiliary perimeter problem}
\label{ct:perimeter-subsection}

Each component of $S$ is a closed orientable surface and bounds a compact
handlebody. Attach such handlebodies along $S$, with the appropriate
orientations, and denote their union by $F$ and the resulting boundaryless
manifold by $\widehat\Omega$. Extend $g$ smoothly across $S$: extend its
smooth coefficients in boundary collars, retain positive definiteness
in smaller collars, and join to any interior metric using a partition
of unity. Thus
\[
 \widehat\Omega\setminus\operatorname{int}F=\Omega,\qquad \partial F=S.
\]
This construction extends only the Riemannian metric.

For a bounded set $M$ of finite perimeter, let
$P_g(M)=|D\mathbf1_M|_g(\widehat\Omega)$. Sets are initially identified
up to volume-null sets. Put
\begin{equation}
\label{ct:perimeter-infimum}
 p_g(F)=\inf\{P_g(M):M\text{ bounded, of finite perimeter, }
                         F\subset M\text{ a.e.}\}.
\end{equation}
Perimeter is taken in all of $\widehat\Omega$; in particular
$P_g(F)=\operatorname{Area}_g(S)$.
The same construction applies to any smooth compact enlarged obstacle
whose exterior is connected and has the same end.

\begin{lemma}[Confinement by the end foliation]
\label{ct:confinement}
There are compact truncations
$C_{R_0}\subset\operatorname{int}C_{R_1}$ containing $F$ such that the
infimum in Equation~\ref{ct:perimeter-infimum} is attained and every
bounded minimizer is contained in $C_{R_0}$ up to a null set. The same
truncations work for a family with uniform asymptotic bounds
$g_{ij}-\delta_{ij}=O_1(r^{-q})$, $q>0$, on a common end.
The existence and confinement statements also hold for an arbitrary
bounded measurable obstacle in place of $F$.
\end{lemma}

\begin{proof}
Let $C_t$ be the compact region inside $S_t$, including the filling.
On a sufficiently distant end the smooth unit field
$Z=\nabla r/|\nabla r|_g$ satisfies
\begin{equation}
\label{ct:positive-divergence}
 \operatorname{div}_g Z=H_{S_r}=\frac2r+O(r^{-1-q})>0.
\end{equation}
Increase $R_0$ so this holds for $r>R_0$ and
$F\subset\operatorname{int}C_{R_0}$.

For any bounded finite-perimeter set $M$, the slicing and Gauss--Green
formulas give, for almost every $t>R_0$,
\begin{align}
 \int_{M\cap\{r>t\}}\operatorname{div}_g Z\,dV_g
 &=\int_{\partial^*M\cap\{r>t\}}g(Z,\nu_M)\,dA_g-b(t),
 \label{ct:clipping-flux}\\
 P_g(M\cap C_t)&=P_g(M;\{r<t\})+b(t),
 \label{ct:clipping-perimeter}
\end{align}
where $b(t)=\int_{S_t}\mathbf1_M\,dA_g$ uses the common
almost-everywhere trace. In Equation~\ref{ct:clipping-flux}, cut off
$Z$ beyond the bounded set $M$ if necessary. Since $|Z|_g=1$,
\begin{equation}
\label{ct:clipping}
 P_g(M\cap C_t)\leq P_g(M)-
       \int_{M\cap\{r>t\}}\operatorname{div}_g Z\,dV_g.
\end{equation}
Clipping strictly decreases perimeter if the discarded volume is
positive.

Fix $R_1>R_0$ and clip each member of a minimizing sequence at a good
radius in $(R_0,R_1)$. The resulting sets lie in $C_{R_1}$, contain $F$,
and have bounded perimeter and volume. Compactness in the space $BV$ of
functions of bounded variation gives an
$L^1$ convergent subsequence, and lower semicontinuity gives a minimizer.
The obstacle condition is closed under $L^1$ convergence. These
compactness and lower-semicontinuity results apply in finitely many
coordinate charts on a larger compact neighborhood
\cite[Theorem~3.23 and Proposition~3.6]{AmbrosioFuscoPallara2000}.
Equation~\ref{ct:clipping} equates this confined infimum with the
unrestricted bounded-set infimum. Applying it to a minimizer at good
radii decreasing to $R_0$ proves that the minimizer has zero volume
outside $C_{R_0}$. Its perimeter support therefore lies strictly inside
$C_{R_1}$.

Uniform end bounds make the error in
Equation~\ref{ct:positive-divergence} uniform, so the same construction
works for the stated family.
Only boundedness and the closedness in $L^1$ of the containment
condition were used for the obstacle; a large compact truncation is
always a finite-perimeter competitor. This proves the last assertion.
\end{proof}

\begin{proposition}[Equivalence with full cuts and obstacle regularity]
\label{ct:perimeter-equivalence}
The auxiliary minimum equals the intrinsic smooth-cut infimum:
\begin{equation}
\label{ct:infima-equal}
 p_g(F)=a_g(S).
\end{equation}
Every minimizing set has a regular representative with compact embedded
$C^{1,1}$ frontier, smooth and minimal away from the obstacle. Thus its
frontier is $C^{1,\alpha}$ for every $0<\alpha<1$ and has $W^{2,2}$
local graphs. It has no bounded complementary component. Its exterior
is connected to the unique end, and its entire frontier, including
contact with $S$, carries its perimeter measure.
\end{proposition}

\begin{proof}
Let $M$ minimize and let $M'$ be a bounded finite-perimeter local
competitor, without an obstacle restriction. The set $M'\cup F$ is
admissible. Perimeter submodularity gives
\begin{equation}
\label{ct:submodular-repair}
 P_g(M)-P_g(M')\leq P_g(F)-P_g(M'\cap F).
\end{equation}
Choose a smooth compactly supported vector field $V$ with $|V|_g\leq1$
which agrees with the outward unit normal of $F$ on $\partial F$.
The signed-distance collar and a cutoff provide such a field.
Gauss--Green gives
\[
 P_g(F)=\int_F\operatorname{div}_gV\,dV_g,\qquad
 \int_{M'\cap F}\operatorname{div}_gV\,dV_g\leq P_g(M'\cap F).
\]
Consequently, for $\Lambda=\|\operatorname{div}_gV\|_\infty$,
\begin{equation}
\label{ct:almost-minimality}
 P_g(M)-P_g(M')\leq
 \Lambda\operatorname{Vol}_g(F\setminus M')
 \leq\Lambda\operatorname{Vol}_g(M\mathbin\triangle M').
\end{equation}
The unchanged perimeters outside the variation neighborhood cancel,
so this is also a local inequality.

The exact obstacle regularity result we use is
Huisken--Ilmanen's Regularity Theorem~1.3(iii): in a smooth Riemannian
domain of ambient dimension less than eight, a pure perimeter minimizer
respecting a $C^2$ obstacle has $C^{1,1}$ boundary, smooth off contact
\cite{HuiskenIlmanen2001}. Here $M$ minimizes against every compactly
supported competitor containing the open obstacle $\operatorname{int}F$;
the metric and obstacle are smooth, the dimension is three, and the
bulk term is zero. Lemma~\ref{ct:confinement} keeps its frontier away
from the auxiliary truncation boundary. The regular representative is
the closure of its measure-theoretic interior, whose boundary equals
the perimeter support. Away from $F$, both signs of every compact
normal variation are allowed, so this smooth boundary has mean
curvature zero. The stated Sobolev regularity follows from $C^{1,1}$.

A compact embedded frontier has finitely many components: it has a
finite cover by connected graph neighborhoods. If its complement had
a bounded open component, that component's boundary would be a nonempty
union of frontier components. Adding it to $M$ removes their positive
area without adding a new frontier and contradicts minimality.
There is only one unbounded complementary component, because
$\widehat\Omega$ has one end and $M$ is bounded.

Given an admissible smooth exterior $D$ with full boundary $\Gamma$,
the bounded set
$M_D=\widehat\Omega\setminus\operatorname{int}D$ contains $F$ and has
boundary exactly $\Gamma$. Therefore
$P_g(M_D)=\operatorname{Area}_g(\Gamma)$, including every coincident
portion on $S$. In particular $D=\Omega$ gives $M_D=F$.
This proves $p_g(F)\leq a_g(S)$.

Conversely, flow a minimizing set $M$ for small time $s>0$ by a
compactly supported smooth field pointing strictly outward on $S$.
The image $M_s$ contains $F$ in its interior. Indeed it contains the
flowed obstacle, whose boundary in the signed-distance collar has
moved strictly outside $S$. The flow's surface Jacobians converge
uniformly to one, so $P_g(M_s)\to P_g(M)$.

A compact $C^1$ embedded boundary has a $C^1$ defining function with
nonzero differential near its boundary. Approximate this function in
$C^1$ by smooth functions; the implicit function theorem yields smooth
embedded boundaries converging in $C^1$, with the same enclosed side.
Apply this to $\partial M_s$. Its positive distance from $F$ preserves
enclosure under sufficiently close approximation, and surface
Jacobians give area convergence. Fill bounded complementary components
of each approximation. This deletes components and does not increase
area. The closure of the remaining exterior is connected, contains
the whole distant end, and has smooth compact two-sided intrinsic
boundary in $\operatorname{int}\Omega$. It is an admissible full cut.
Letting the approximation error and then $s$ tend to zero proves
$a_g(S)\leq p_g(F)$.
\end{proof}

\begin{corollary}[Minimal obstacles]
\label{ct:minimal-obstacle}
If every component of the smooth obstacle boundary is minimal, every
minimizing frontier in Proposition~\ref{ct:perimeter-equivalence} is
smooth and minimal everywhere. If a frontier meets a connected
component of the obstacle boundary, it contains that entire component.
\end{corollary}

\begin{proof}
At contact, write frontier and obstacle as graphs $u\geq\psi$, with
$u\in C^{1,1}$ and $\psi$ smooth. On the contact set $Du=D\psi$ and
$D^2u=D^2\psi$ almost everywhere. The second statement follows from
the fact that the weak derivative of a Sobolev function vanishes
almost everywhere on a level set, applied to the Lipschitz functions
$D_i(u-\psi)$. Off contact $u$ solves the minimal graph equation; at
almost every contact point its second derivatives equal those of
the minimal graph $\psi$. Thus it solves the same equation weakly
throughout the patch. The bounded gradient makes this a uniformly
elliptic equation, so interior elliptic regularity and bootstrapping
give smoothness. The strong maximum principle for the difference of
the two smooth minimal graph equations gives local coincidence at
each contact point. Contact is also closed, so connectedness extends
the coincidence over the obstacle component.
\end{proof}

\begin{corollary}[Comparison of exterior metrics]
\label{ct:metric-comparison}
If $g_j\geq c_jg$ on $\Omega$, $c_j\to1$, and $g_j\to g$ uniformly
on compact sets, then $a_{g_j}(S)\to a_g(S)$.
\end{corollary}

\begin{proof}
Two-dimensional area elements satisfy $dA_{g_j}\geq c_jdA_g$, so
$a_{g_j}(S)\geq c_ja_g(S)$. For the upper bound, test with any fixed
smooth cut of $g$-area less than $a_g(S)+\varepsilon$, use compact
convergence on that cut, and then let $\varepsilon\downarrow0$.
\end{proof}

\subsection{Strict compactness and the derivative of the minimum}
\label{ct:variation-subsection}

For a fixed smooth obstacle $F$, denote the family of its minimizing
filled sets by $\mathcal M_g(F)$. Give this family the $L^1$ distance
on a common compact truncation; enlarging that truncation does not
change the distance. Tangent planes are unoriented.

\begin{proposition}[Strict compactness of minimizing enclosures]
\label{ct:strict-compactness}
Let smooth metrics $g_j$ converge to $g$ uniformly on compact sets,
with common end bounds as in Lemma~\ref{ct:confinement}, and choose
their auxiliary extensions to converge uniformly on the filling.
If $M_j\in\mathcal M_{g_j}(F)$, a subsequence converges to
$M\in\mathcal M_g(F)$ with
\begin{equation}
\label{ct:strict-convergence}
 \mathbf1_{M_j}\longrightarrow\mathbf1_M\quad\text{in }L^1,\qquad
 P_g(M_j)\longrightarrow P_g(M)=p_g(F).
\end{equation}
For every continuous function $\Psi$ on the bundle of unoriented
two-planes over the common compact truncation,
\begin{equation}
\label{ct:plane-measure-continuity}
 \int_{\partial M_j}\Psi(x,T_x\partial M_j)\,dA_g
 \longrightarrow\int_{\partial M}\Psi(x,T_x\partial M)\,dA_g.
\end{equation}
Thus $\mathcal M_g(F)$ is a compact metric space, and the area
first-variation functional is continuous on it.
\end{proposition}

\begin{proof}
Confinement supplies one compact set $C$ containing every minimizer.
Uniform metric comparison on $C$ gives $\varepsilon_j\downarrow0$
such that for every set $M'$ whose perimeter is supported there,
\begin{equation}
\label{ct:uniform-perimeter-comparison}
 (1-\varepsilon_j)P_g(M')\leq P_{g_j}(M')
 \leq(1+\varepsilon_j)P_g(M').
\end{equation}
Taking infima, using confinement for both metrics, proves
$p_{g_j}(F)\to p_g(F)$, and hence $P_g(M_j)\to p_g(F)$.
Compactness in $BV$ and lower semicontinuity give a subsequential
$L^1$ limit $M$ containing $F$ and supported in $C$, with
$P_g(M)\leq p_g(F)$. Therefore it minimizes, proving
Equation~\ref{ct:strict-convergence}.

The vector-measure continuity needed here can also be proved directly.
Use the fixed metric $g$, and write
\[
 \lambda_j=\nu_j\,dA_g|_{\partial M_j},\qquad
 \lambda=\nu\,dA_g|_{\partial M},
\]
where $\nu_j,\nu$ are outward unit normals. Gauss--Green and $L^1$
convergence give weak convergence of these vector measures.
Their total variation masses converge by
Equation~\ref{ct:strict-convergence}. Lower semicontinuity with
nonnegative continuous weights and convergence of the total masses
also imply $|\lambda_j|\rightharpoonup|\lambda|$.

Approximate the measurable unit normal $\nu$ in $L^2(|\lambda|)$ by
a continuous vector field $V$ on $C$ with $|V|_g\leq1$. One obtains
such approximations in finitely many bundle charts, joins them by a
partition of unity, and projects onto the unit ball. They make
$\int(1-g(V,\nu))\,d|\lambda|$ arbitrarily small. For each such $V$,
\begin{align*}
 \int|\nu_j-V|_g^2\,d|\lambda_j|
 &\leq2\left(|\lambda_j|(C)-\int g(V,d\lambda_j)\right),\\
 \limsup_j\int|\nu_j-V|_g^2\,d|\lambda_j|
 &\leq2\int(1-g(V,\nu))\,d|\lambda|.
\end{align*}
A continuous function of a unit normal extends continuously to the
closed unit-ball bundle and is uniformly continuous on that compact
bundle. The displayed estimate, Cauchy--Schwarz, and weak convergence
of $|\lambda_j|$ give convergence of its integrals, first at a fixed
approximation $V$, then as the approximation error tends to zero.
Apply this to the even function $\Psi(x,\nu^{\perp_g})$ to obtain
Equation~\ref{ct:plane-measure-continuity}. This is the needed case of
Reshetnyak's Continuity Theorem
\cite[Theorem~2.39]{AmbrosioFuscoPallara2000}.

The constant metric sequence gives compactness of $\mathcal M_g(F)$.
For a continuous symmetric tensor $h$, the function
$\Psi(x,T)=\tfrac12\operatorname{tr}_T h$ is continuous, giving the
last assertion.
\end{proof}

\begin{proposition}[One-sided derivative of the area infimum]
\label{ct:area-derivative}
Let $g_s$ be a path of smooth exterior metrics with $g_0=g$, common
end bounds as in Lemma~\ref{ct:confinement}, and
\[
 g_s=g+sh+o(|s|)\quad\text{in }C^0
\]
on a common compact truncation including the filling. For the fixed
obstacle $F$ with exterior boundary $B$, set $a(s)=a_{g_s}(B)$. Then
\begin{equation}
\label{ct:minimum-derivative}
 \left.\frac{d}{ds}\right|_{0+}a(s)
 =\min_{M\in\mathcal M_g(F)}L_h(M),\qquad
 L_h(M)=\frac12\int_{\partial M}
                  \operatorname{tr}_{T\partial M}h\,dA_g.
\end{equation}
Only exterior and boundary values of $h$ enter the integral.
\end{proposition}

\begin{proof}
The determinant formula for a two-plane area gives, uniformly over
planes in the compact truncation,
\[
 dA_{g_s}=\left(1+\frac s2\operatorname{tr}_T h+o(|s|)\right)dA_g.
\]
For every finite-perimeter competitor $M'$ in that truncation,
\begin{equation}
\label{ct:uniform-area-expansion}
 P_{g_s}(M')=P_g(M')+sL_h(M')+o(|s|)P_g(M'),
\end{equation}
with the same remainder modulus. By
Proposition~\ref{ct:strict-compactness}, the minimum of $L_h$ exists.

Testing $a(s)$ with any $M\in\mathcal M_g(F)$ gives
\[
 \limsup_{s\downarrow0}\frac{a(s)-a(0)}s\leq L_h(M).
\]
In the other direction, take a minimizer $M_s$ for $g_s$. Its
$g$-perimeter is bounded and at least $a(0)$, so
Equation~\ref{ct:uniform-area-expansion} gives
\[
 \frac{a(s)-a(0)}s\geq L_h(M_s)+o(1)\qquad(s>0).
\]
Along a sequence realizing the lower limit, strict compactness gives
a subsequence converging to a member of $\mathcal M_g(F)$; continuity
of $L_h$ yields the opposite inequality. No differentiable choice of
minimizer is needed. Since $F\subset M$, the frontier does not enter
$\operatorname{int}F$, which proves the final assertion.
\end{proof}

\subsection{Compatibility of minimizing sheets}
\label{ct:sheets-subsection}

The variation argument averages over the minimizing family. The next
property prevents cancellation between different tangent planes above
one spatial point.

\begin{lemma}[Compatible planes]
\label{ct:compatible-planes}
For a fixed smooth compact obstacle $F$, the set
\[
 \mathcal Z=\bigcup_{M\in\mathcal M_g(F)}\partial M
\]
is compact. At each $x\in\mathcal Z$, all minimizing frontiers through
$x$ have the same unoriented tangent plane $T(x)$, and $T$ is continuous
on $\mathcal Z$. If $\pi$ is a Borel probability measure on
$\mathcal M_g(F)$ and
\[
 \eta(A)=\int_{\mathcal M_g(F)}
          \operatorname{Area}_g(A\cap\partial M)\,d\pi(M),
\]
then, for every continuous function $\Psi$ on the plane bundle,
\begin{equation}
\label{ct:averaged-plane}
 \int_{\mathcal M_g(F)}\int_{\partial M}
        \Psi(x,T_x\partial M)\,dA_g\,d\pi(M)
 =\int_{\mathcal Z}\Psi(x,T(x))\,d\eta(x).
\end{equation}
\end{lemma}

\begin{proof}
If $M_1,M_2$ minimize, their union and intersection contain $F$.
Submodularity gives
\[
 2p_g(F)\leq P_g(M_1\cup M_2)+P_g(M_1\cap M_2)
 \leq P_g(M_1)+P_g(M_2)=2p_g(F).
\]
Thus both also minimize. If their frontier planes were distinct at
an intersection, the tangent half-spaces would meet transversely.
The boundary of their union or intersection would have a corner:
its tangent cone would contain two distinct half-planes. This
contradicts its $C^1$ regularity from
Proposition~\ref{ct:perimeter-equivalence}.

Here is the required closedness argument. On the common compact
truncation, Equation~\ref{ct:almost-minimality} has a fixed constant.
It yields uniform density bounds at every frontier point, for
sufficiently small coordinate balls centered there:
\begin{equation}
\label{ct:density-bounds}
 \begin{gathered}
 c r^3\leq\operatorname{Vol}_g(M\cap B_r),\qquad
 c r^3\leq\operatorname{Vol}_g(B_r\setminus M),\\
 c r^2\leq P_g(M;B_r)\leq C r^2.
 \end{gathered}
\end{equation}
For clarity, compare with $M\setminus B_r$ and $M\cup B_r$.
For either phase volume $v(r)$, the isoperimetric inequality for its
intersection with the ball and these comparisons give
\[
 v(r)^{2/3}\leq C\bigl(2v'(r)+\Lambda v(r)\bigr)
 \quad\text{for almost every }r.
\]
The coarea formula identifies $v'$ with the slice area up to fixed
coordinate constants. Since $v(r)\leq Cr^3$, the last term is absorbed
for small $r$. Integrating $(v^{1/3})'\geq c$ gives $v(r)\geq cr^3$;
both phases are nonempty at every scale at a perimeter-support point.
Relative isoperimetry gives the perimeter lower bound; the same
competitors and the coordinate-sphere area bound give the upper bound.
Smooth compact geometry makes these constants uniform.

Suppose now $M_j\to M$ in $\mathcal M_g(F)$ and
$x_j\in\partial M_j$ tend to $x$. The two volume bounds pass to every
fixed small ball about $x$ under $L^1$ convergence; hence
$x\in\partial M$. Huisken--Ilmanen's Regularity Theorem~1.3(ii)
gives uniform local $C^{1,\alpha}$ estimates for a fixed
$0<\alpha\leq1/2$ \cite{HuiskenIlmanen2001}. Its dependencies are
the distance to the ambient domain boundary, obstacle
$C^{1,\alpha}$ bounds, metric $C^1$ bounds and positive definiteness,
and the bulk bound, here zero. All are fixed in this application.
Arzel\`a--Ascoli in local graph charts gives subsequential
$C^{1,\beta}$ convergence, $0<\beta<\alpha$. The density bounds and
the $L^1$ limit identify the limiting sheet with $\partial M$. Thus
\[
 T_{x_j}\partial M_j\longrightarrow T_x\partial M.
\]
The incidence set
$\{(M,x,T_x\partial M):M\in\mathcal M_g(F),\,x\in\partial M\}$
is therefore compact. Its spatial projection $\mathcal Z$ is compact,
and the previously proved uniqueness of the plane at each limit point
makes $T$ continuous.

By Proposition~\ref{ct:strict-compactness}, the surface measures
depend weakly continuously on $M$, so their average defines the
finite Borel measure $\eta$. Each surface uses the same plane $T(x)$
at its points. Integrating this identity proves
Equation~\ref{ct:averaged-plane}.
\end{proof}
